\documentclass[11pt]{article}
\usepackage[margin=1in]{geometry}
\usepackage{amsmath,amssymb,amsthm,hyperref}
\hypersetup{colorlinks=true,urlcolor=blue,linkcolor=blue}
\newtheorem{theorem}{Theorem}
\newtheorem{lemma}{Lemma}
\newtheorem{corollary}{Corollary}
\newcommand{\N}{\mathbb N}
\title{Odd moduli refute eventual coefficient monotonicity\\
\large MathDB 344349}
\author{Principia Math In Progress}
\date{18 September 2026}
\begin{document}
\maketitle

\section*{Task 1: full resolution of the stated conjecture}
Put
\[
 F_m(q)=\frac1{(q,-q^{m-1};q^m)_\infty}
       =\sum_{n\ge0}a_m(n)q^n,
 \qquad (u;v)_\infty=\prod_{j\ge0}(1-uv^j),
\]
and put $a_m(n)=0$ for $n<0$.
Seo and Yee conjectured that for every $m\ge4$ there is an $N(m)$
such that $a_m(n)\ge a_m(n-m)$ for every $n\ge N(m)$
\cite[Conjecture 22]{SY}. This is Conjecture 5 in their arXiv version.
We disprove it within its stated range.

\begin{theorem}\label{thm:main}
For every odd integer $m\ge5$, there are infinitely many odd integers
$n$ such that
\[
                     a_m(n)<a_m(n-m).
\]
Consequently no threshold $N(m)$ exists for any such $m$.
\end{theorem}

All products and power series below converge absolutely on compact
subsets of $|q|<1$. No sign assumption on the coefficients is made.
For positive functions, $U(t)=\Theta_m(V(t))$ means that their ratio
is bounded above and below by positive constants depending on $m$,
for all sufficiently small positive $t$.

\begin{lemma}\label{lem:ratio}
For every fixed odd integer $m\ge3$ and $r=e^{-t}$ with $t\downarrow0$,
\[
                   \frac{F_m(-r)}{F_m(r)}
                         =\Theta_m(t^{1-1/m}).
\]
\end{lemma}
\begin{proof}
Split the product indices into even and odd indices and set $Q=r^{2m}$.
Since $m$ is odd, this gives
\begin{align*}
 F_m(r)&=\frac1{(r,r^{m+1},-r^{m-1},-r^{2m-1};Q)_\infty},\\
 F_m(-r)&=\frac1{(-r,r^{m+1},-r^{m-1},r^{2m-1};Q)_\infty}.
\end{align*}
Their ratio $R_m(r)$ is positive and equals
\begin{equation}\label{eq:product}
 R_m(r)=\prod_{j\ge0}
 \frac{1-r^{2mj+1}}{1+r^{2mj+1}}
 \frac{1+r^{2mj+2m-1}}{1-r^{2mj+2m-1}}.
\end{equation}
The absolutely convergent expansion
$\log(1-u)-\log(1+u)=-2\sum_{k\ge1,\ k\text{ odd}}u^k/k$
therefore yields
\begin{equation}\label{eq:log}
 \log R_m(e^{-t})
 =-2\sum_{\substack{k\ge1\\k\text{ odd}}}\frac{f_m(kt)}{k},
 \qquad
 f_m(x)=\frac{e^{-x}-e^{-(2m-1)x}}{1-e^{-2mx}}
       =\frac{\sinh((m-1)x)}{\sinh(mx)}.
\end{equation}
Let $c=(m-1)/m$. Taylor expansion at zero and the exponential
expression above give constants $C_m$ such that
\[
 |f_m(x)-c|\le C_m x^2\quad(0<x\le1),\qquad
 0<f_m(x)\le C_m e^{-x}\quad(x\ge1).
\]
It follows that
\[
 \sum_{\substack{k\le1/t\\k\text{ odd}}}
       \frac{f_m(kt)-c}{k}=O_m(1),
 \qquad
 \sum_{\substack{k>1/t\\k\text{ odd}}}
       \frac{f_m(kt)}{k}=O_m(1).
\]
Indeed the absolute values are bounded respectively by
$C_m t^2\sum_{k\le1/t}k$ and
$C_m t\sum_{k>1/t}e^{-kt}$, both bounded as $t\downarrow0$.
The elementary odd harmonic sum satisfies
\[
 \sum_{\substack{k\le X\\k\text{ odd}}}\frac1k
                 =\frac12\log X+O(1).
\]
Substitution into \eqref{eq:log} now gives
$\log R_m(e^{-t})=c\log t+O_m(1)$.
Exponentiation proves the lemma, including both positive bounds.
\end{proof}

\begin{proof}[Proof of Theorem~\ref{thm:main}]
Write
\[
 D_m(q)=(1-q^m)F_m(q)=\sum_{n\ge0}d_m(n)q^n,
 \qquad d_m(n)=a_m(n)-a_m(n-m).
\]
For $0<r<1$, both $D_m(r)$ and $D_m(-r)$ are positive, and the lemma gives
\begin{equation}\label{eq:dominance}
 \frac{D_m(-r)}{D_m(r)}
 =\frac{1+r^m}{1-r^m}R_m(r)
 =\Theta_m(t^{-1/m})\longrightarrow+\infty,
 \qquad r=e^{-t}.
\end{equation}
We also need $D_m(r)\to+\infty$. For every $j\ge0$,
$r^{mj+m-1}\le r^{mj+1}$, so
\[
 (1-r^{mj+1})(1+r^{mj+m-1})\le1-r^{2mj+2}.
\]
Each reciprocal factor is at least one. Keeping the first two
factors therefore proves
\[
 D_m(r)\ge
 \frac{1-r^m}{(1-r^2)(1-r^{2m+2})}
 \sim \frac{m}{4(m+1)(1-r)}\longrightarrow+\infty.
\]
Together with \eqref{eq:dominance}, this shows that
$D_m(-r)-D_m(r)\to+\infty$. But absolute convergence gives
\[
 D_m(-r)-D_m(r)
                =-2\sum_{\substack{n\ge0\\n\text{ odd}}}d_m(n)r^n.
\]
If only finitely many odd $n$ had $d_m(n)<0$, the expression on the
right would be bounded above, uniformly in $0<r<1$, by twice the
sum of the absolute values of those finitely many negative coefficients.
This contradiction proves the theorem.
\end{proof}

\begin{corollary}
For every odd $m\ge3$ and every positive odd step $h$, infinitely
many odd $n$ satisfy $a_m(n)<a_m(n-h)$.
\end{corollary}
\begin{proof}
Use $(1-q^h)F_m(q)$ in place of $D_m(q)$. The factor in
\eqref{eq:dominance} becomes $(1+r^h)/(1-r^h)\sim2/(ht)$.
The same two-factor lower bound, with numerator $1-r^h$, still
diverges. The odd-coefficient argument is unchanged.
\end{proof}

\paragraph{Verification and scope.}
The accompanying exact verifier independently generates coefficients
by product division and by a logarithmic-derivative recurrence,
checks the product identity \eqref{eq:product} as a formal series,
and checks its logarithmic identity. These are finite consistency
checks; the infinite conclusion follows from the displayed estimates
and the boundedness contradiction. This proof does not assert eventual
negativity at every odd index, or decide even moduli or even steps.
The separate coefficient-nonnegativity conjecture is also not decided.

\section*{Four proposed follow-up tasks}
\begin{description}
\item[Task 2.] Determine whether $a_m(n)\ge a_m(n-m)$ holds for all
sufficiently large $n$ for every even $m\ge4$.
\item[Task 3.] For each odd $m\ge5$, determine whether
$a_m(n)<a_m(n-m)$ holds at every sufficiently large odd $n$;
obtain an effective threshold if it does.
\item[Task 4.] Prove or refute eventual monotonicity with step $2m$
for every odd $m\ge5$.
\item[Task 5.] For coprime integers $1\le a<m/2$, classify the positive
steps $h$ for which the coefficients of
$(q^a,-q^{m-a};q^m)_\infty^{-1}$ are eventually nondecreasing
under a shift by $h$.
\end{description}
These are proposed research tasks and earn no additional solution credit.

\begin{thebibliography}{9}
\bibitem{SY} S. Seo and A. J. Yee,
\emph{Index of seaweed algebras and integer partitions},
Electronic Journal of Combinatorics \textbf{27}(1) (2020), P1.47,
\href{https://doi.org/10.37236/9054}{doi:10.37236/9054}.
Preprint: \href{https://arxiv.org/abs/1910.14369}{arXiv:1910.14369}.
\end{thebibliography}
\end{document}
